% ==============================================================================
% CAPÍTULO 3: SEGURANÇA EM BANCO DE DADOS: CONTROLE DE ACESSO, PRIVILÉGIOS E ROLES
% ==============================================================================

\begin{center}
\setstretch{1.0}
\textbf{\fontsize{14}{16}\selectfont CAPÍTULO 3}

\vspace{0.5em}

\textbf{\fontsize{14}{16}\selectfont SEGURANÇA EM BANCO DE DADOS: CONTROLE DE ACESSO, MATRIZES DE PRIVILÉGIOS E RBAC}
\end{center}

\vspace{1.5em}

\section*{1. CONTEXTUALIZAÇÃO DO TEMA E A IMPORTÂNCIA DO APRENDIZADO}

Nos ecossistemas corporativos contemporâneos, o banco de dados representa o cofre digital da organização, concentrando registros financeiros, estratégias de mercado, segredos industriais e milhões de dados pessoais de clientes e cidadãos. Consequentemente, os SGBDs são o alvo prioritário de agentes maliciosos, ataques cibernéticos patrocinados e fraudes internas. Nesse cenário, o domínio formal da \textbf{Segurança em Bancos de Dados} e dos \textbf{Modelos de Controle de Acesso} transcende a mera rotina operacional: é um imperativo de sobrevivência e conformidade institucional.

Estudar segurança em nível de banco de dados é indispensável porque a segurança perimetral (como firewalls de rede e criptografia de canal) não protege a organização contra ameaças originadas em credenciais comprometidas, falhas em códigos de aplicação ou acessos indevidos por colaboradores com privilégios excessivos.

\subsection*{1.1. A Problemática Crítica do Controle de Acesso Inadequado}

A negligência na segregação de acessos e permissões no SGBD desencadeia vulnerabilidades estruturais recorrentes no mercado de trabalho:
\begin{itemize}
    \item \textbf{O Vício Corporativo da Conta de Superusuário Compartilhada:} Uma das más práticas mais nocivas da indústria é conectar aplicações web, microsserviços, desenvolvedores e estagiários utilizando a mesma credencial de superusuário (\texttt{root} ou \texttt{sa}). Sob esse arranjo, qualquer falha em uma tela da aplicação concede controle irrestrito sobre todo o banco de dados, permitindo a exclusão de esquemas inteiros e a visualização de dados sigilosos.
    \item \textbf{Ataques por Injeção de SQL (\textit{SQL Injection}):} Listado consistentemente entre os riscos mais críticos do guia internacional \textit{OWASP Top 10} (OWASP, 2021), o ataque de injeção permite que atacantes externos subvertam consultas e burlem telas de login quando entradas de usuários são concatenadas diretamente em código SQL sem parametrização.
    \item \textbf{Violação do Princípio do Menor Privilégio:} Quando não há segregação de papéis, funcionários do setor de atendimento podem acessar salários de executivos, professores podem alterar tabelas de configuração ou alunos podem visualizar notas de terceiros. A falta de granularidade por coluna e por tabela expõe a instituição a fraudes e processos judiciais.
    \item \textbf{Proliferação e Descontrole de Permissões:} Em ambientes que não utilizam o modelo de papéis (\textit{Role-Based Access Control} -- RBAC), permissões são concedidas individualmente a cada funcionário. Quando um colaborador muda de setor ou é desligado da empresa, suas permissões antigas permanecem ativas na base de dados por anos a fio, criando portas dos fundos desmonitoradas.
\end{itemize}

O aprendizado deste capítulo instrumentaliza o leitor a implementar o célebre \textbf{Princípio do Menor Privilégio} (\textit{Principle of Least Privilege}), assegurando que cada usuário, aplicação ou serviço receba estritamente as permissões mínimas e necessárias para desempenhar sua função funcional, e nada além disso.

\newpage

\section*{2. FUNDAMENTAÇÃO TEÓRICA DA SEGURANÇA EM SGBDs}

Os alicerces formais da segurança computacional aplicada a sistemas de informação foram estabelecidos por Saltzer e Schroeder (1975) em seu influente trabalho sobre proteção de informações em computadores. Dentre os princípios formulados, destacam-se a economia de mecanismos, a separação de privilégios e a exigência de que o acesso padrão a qualquer recurso seja rigorosamente nulo até que uma concessão explícita seja homologada.

A literatura moderna de segurança da informação (STALLINGS; BROWN, 2018; SILBERSCHATZ; KORTH; SUDARSHAN, 2020) decompõe a segurança no SGBD em duas barreiras distintas e complementares:
\begin{enumerate}
    \item \textbf{Autenticação:} O mecanismo pelo qual o sistema verifica a veracidade da identidade alegada por um agente externo (``Quem é você?''). Em bancos de dados relacionais, isso é realizado pela combinação de usuário, máquina/host de origem (\texttt{'usuario'@'host'}) e segredos criptográficos ou certificados.
    \item \textbf{Autorização:} O procedimento contínuo pelo qual o motor do SGBD consulta o catálogo de privilégios para determinar se uma identidade autenticada possui o direito legítimo de executar uma instrução DDL ou DML sobre um determinado objeto (``O que você tem permissão de executar?'').
\end{enumerate}

Para estruturar a autorização, a ciência da computação desenvolveu três modelos clássicos de controle de acesso (SANDHU et al., 1996; ELMASRI; NAVATHE, 2011):
\begin{itemize}
    \item \textbf{DAC (\textit{Discretionary Access Control}):} Controle discricionário, onde o criador/dono do objeto decide quem recebe privilégios.
    \item \textbf{MAC (\textit{Mandatory Access Control}):} Controle obrigatório, baseado em políticas impositivas e níveis de classificação de segurança.
    \item \textbf{RBAC (\textit{Role-Based Access Control}):} Controle baseado em papéis organizacionais, o padrão ouro das arquiteturas corporativas modernas.
\end{itemize}

\newpage

\section*{3. INJEÇÃO DE SQL (SQL INJECTION): ANATOMIA E DEFESA}

O ataque de \textbf{Injeção de SQL} ocorre quando entradas textuais fornecidas por usuários são concatenadas de forma direta em instruções SQL dinâmicas sem sanitização ou parametrização, permitindo que a entrada do atacante altere a sintaxe e a lógica da consulta planejada pelo desenvolvedor (SILBERSCHATZ; KORTH; SUDARSHAN, 2020).

\subsection*{3.1. Demonstração Prática da Vulnerabilidade}

Considere uma relação de credenciais de usuários, na qual as senhas são armazenadas como resumo criptográfico (hash SHA-256):

\begin{lstlisting}
USE escola;

CREATE TABLE usuario (
    id INT PRIMARY KEY AUTO_INCREMENT,
    login VARCHAR(30) UNIQUE NOT NULL,
    senha_hash CHAR(64) NOT NULL
);

-- Insercao com hash criptografico SHA-256
INSERT INTO usuario (login, senha_hash) VALUES
('ana', SHA2('ana#123', 256)),
('admin', SHA2('Adm!n#2026', 256));
\end{lstlisting}

Se a aplicação web autenticar os usuários concatenando strings da seguinte forma:
\begin{lstlisting}
SELECT * FROM usuario WHERE login = 'ana';
\end{lstlisting}

Um invasor pode submeter no campo de login a carga útil (\textit{payload}) clássica: \texttt{' OR '1'='1}. A consulta processada pelo motor do SGBD resultará em:

\begin{lstlisting}
SELECT * FROM usuario WHERE login = '' OR '1'='1';
\end{lstlisting}
Como a condição booleana \texttt{'1'='1'} é invariavelmente verdadeira (\textit{tautologia}), o SGBD retorna todas as tuplas da tabela, concedendo acesso com os privilégios do primeiro usuário listado (comumente o administrador do sistema).

\subsection*{3.2. A Defesa Efetiva: Consultas Preparadas (Prepared Statements)}

A única mitigação definitiva contra a injeção de SQL consiste na dissociação estrita entre o plano de compilação do código executável e os parâmetros de dados. No MySQL, isso é garantido pelo uso de \textbf{Comandos Preparados} (\textit{Prepared Statements}):

\begin{lstlisting}
-- 1. Compilacao previa da instrucao SQL com marcador de parametro '?'
PREPARE stmt FROM 'SELECT id, login FROM usuario WHERE login = ?';

-- 2. Teste do payload malicioso passado como dado puro
SET @login_injetado = "' OR '1'='1";
EXECUTE stmt USING @login_injetado;
-- O motor busca literalmente um usuario cujo login seja "' OR '1'='1" e retorna zero tuplas!

-- 3. Execucao legitima com parametro autentico
SET @login_valido = 'ana';
EXECUTE stmt USING @login_valido;

-- 4. Desalocacao dos recursos no servidor
DEALLOCATE PREPARE stmt;
\end{lstlisting}

Com a parametrização, os dados inseridos jamais são interpretados como operadores lógicos pelo analisador sintático do banco de dados.

\newpage

\section*{4. GESTÃO DE CONTAS E POLÍTICAS DE SENHAS NO MYSQL}

No MySQL, a identidade de uma conta é composta pelo nome do usuário e pelo computador/rede de origem da conexão (\texttt{'usuario'@'host'}). Isso permite isolar conexões locais (\texttt{'localhost'}) de acessos remotos via rede corporativa (ex.: \texttt{'192.168.1.\%'}).

\subsection*{4.1. Criação de Contas e Restrições de Conexão}

\begin{lstlisting}
-- Criacao de conta restrita ao host local
CREATE USER 'joao'@'localhost' IDENTIFIED BY 'SenhaForte#2025';

-- Consulta ao catalogo de sistema de usuarios
SELECT user, host, plugin FROM mysql.user WHERE user = 'joao';
\end{lstlisting}

\subsection*{4.2. Políticas Corporativas de Senhas e Bloqueio por Força Bruta}

Para atender a padrões corporativos rigorosos (como normas ISO/IEC 27001 e guias CIS Benchmark), o DBA deve impor políticas de ciclo de vida e proteção contra tentativas repetidas de invasão:

\begin{lstlisting}
-- 1. Expiracao periodica obrigatoria a cada 90 dias
ALTER USER 'joao'@'localhost' PASSWORD EXPIRE INTERVAL 90 DAY;

-- 2. Bloqueio temporario da conta apos 3 tentativas consecutivas incorretas
ALTER USER 'joao'@'localhost'
FAILED_LOGIN_ATTEMPTS 3
PASSWORD_LOCK_TIME 1;
\end{lstlisting}

\subsection*{4.3. A Separação entre Autenticar e Autorizar}

Quando o usuário João estabelece sua conexão:
\begin{lstlisting}
SELECT CURRENT_USER(); -- Retorna 'joao@localhost' (Autenticado com sucesso)

-- Tentativa de consulta na tabela de alunos
-- Retorno do motor: Error 1142 (42000): SELECT command denied to user 'joao'@'localhost' for table 'aluno'
SELECT * FROM escola.aluno;
\end{lstlisting}

Em observância ao Princípio do Menor Privilégio, uma conta recém-criada possui privilégios estritamente nulos sobre as tabelas da aplicação até que uma concessão explícita seja homologada.

\newpage

\section*{5. MODELOS CLÁSSICOS DE CONTROLE DE ACESSO}

\subsection*{5.1. Controle de Acesso Discricionário (DAC)}

No modelo DAC (\textit{Discretionary Access Control}), o proprietário de um recurso concede ou revoga permissões aos demais usuários por meio das cláusulas \texttt{GRANT} e \texttt{REVOKE} (ELMASRI; NAVATHE, 2011).

\begin{lstlisting}
-- Concessao granular nos niveis de SELECT e INSERT para o usuario Joao
GRANT SELECT, INSERT ON escola.aluno TO 'joao'@'localhost';

-- Verificacao dos privilegios efetivos concedidos a conta
SHOW GRANTS FOR 'joao'@'localhost';

-- Revogacao pontual do privilegio de gravacao
REVOKE INSERT ON escola.aluno FROM 'joao'@'localhost';
\end{lstlisting}

\subsubsection*{5.1.1. Granularidade por Coluna (Proteção de Campos Confidenciais)}

O SGBD permite conceder acessos restritos a colunas específicas, resguardando atributos sensíveis (como CPF ou salário):

\begin{lstlisting}
CREATE USER 'maria'@'localhost' IDENTIFIED BY 'Maria#2025';

-- Concedendo acesso exclusivamente a colunas operacionais
GRANT SELECT (id_aluno, nome, email) ON escola.aluno TO 'maria'@'localhost';

-- Teste conectado como 'maria':
SELECT nome, email FROM escola.aluno; -- Operacao autorizada

-- Tentativa de acessar a coluna CPF:
-- Retorno do motor: Error 1142 (42000): SELECT command denied on column 'cpf'
SELECT cpf FROM escola.aluno;
\end{lstlisting}

\subsubsection*{5.1.2. O Risco da Cláusula WITH GRANT OPTION}

A cláusula \texttt{WITH GRANT OPTION} permite que o usuário repasse suas próprias permissões a terceiros. Esse poder deve ser monitorado com rigor extremo para evitar a proliferação desgovernada de privilégios administrativos:

\begin{lstlisting}
GRANT SELECT ON escola.aluno TO 'prof_ana'@'localhost' WITH GRANT OPTION;

-- Revogacao exclusiva do poder de repasse mantendo o privilegio original
REVOKE GRANT OPTION ON escola.aluno FROM 'prof_ana'@'localhost';
\end{lstlisting}

\subsection*{5.2. Controle de Acesso Obrigatório (MAC)}

No modelo MAC (\textit{Mandatory Access Control}), os acessos são regidos por uma política central impositiva. Os dados recebem rótulos de segurança (ex.: 1 -- Público, 2 -- Confidencial, 3 -- Sigiloso) e os usuários recebem autorizações de credenciamento (\textit{clearance}) equivalentes (STALLINGS; BROWN, 2018).

No MySQL Community, essa arquitetura é comumente implementada com visões (\textit{views}) que filtram dinamicamente as tuplas cruzando a identidade da sessão ativa (\texttt{USER()}) com a tabela de habilitação:

\begin{lstlisting}
CREATE TABLE documento (
    id INT PRIMARY KEY AUTO_INCREMENT,
    titulo VARCHAR(80) NOT NULL,
    nivel TINYINT NOT NULL CHECK (nivel BETWEEN 1 AND 3)
);

CREATE TABLE habilitacao (
    usuario VARCHAR(30) PRIMARY KEY,
    nivel TINYINT NOT NULL
);

INSERT INTO documento (titulo, nivel) VALUES
('Cardapio Escolar', 1), ('Planilha de Notas', 2), ('Ata Disciplinar Secreta', 3);

INSERT INTO habilitacao VALUES ('prof_ana', 2), ('prof_bia', 1);

-- Visao com filtro dinamico baseado na identidade da conexao
CREATE VIEW documento_permitido AS
SELECT d.id, d.titulo, d.nivel
FROM documento d
JOIN habilitacao h ON h.usuario = SUBSTRING_INDEX(USER(), '@', 1)
WHERE d.nivel <= h.nivel;
\end{lstlisting}

Ao liberar privilégios exclusivamente sobre a visão \texttt{documento\_permitido}, cada usuário visualiza estritamente os documentos compatíveis com o seu nível de credenciamento.

\newpage

\section*{6. CONTROLE DE ACESSO BASEADO EM PAPÉIS (RBAC)}

Em médias e grandes corporações, atribuir privilégios diretamente a usuários individuais é inviável e acarreta severos riscos operacionais. O modelo \textbf{RBAC} (\textit{Role-Based Access Control}) soluciona essa complexidade associando privilégios a \textbf{Papéis} (\textit{Roles}), que representam funções e cargos organizacionais. Os colaboradores recebem um ou mais papéis pertinentes ao seu perfil (SANDHU et al., 1996).

\subsection*{6.1. Ciclo de Vida de Papéis no MySQL}

\begin{lstlisting}
-- 1. Criacao dos papeis institucionais
CREATE ROLE 'r_secretaria', 'r_professor', 'r_aluno';

-- 2. Concessao de privilegios direcionados aos papeis
GRANT SELECT, INSERT, UPDATE, DELETE ON escola.aluno TO 'r_secretaria';
GRANT SELECT ON escola.aluno TO 'r_professor';

-- 3. Criacao do colaborador e atribuicao do papel
CREATE USER 'sec_carla'@'localhost' IDENTIFIED BY 'Carla#2025';
GRANT 'r_secretaria' TO 'sec_carla'@'localhost';

-- 4. Ativacao do papel padrao ao conectar
SET DEFAULT ROLE 'r_secretaria' TO 'sec_carla'@'localhost';
\end{lstlisting}

\subsection*{6.2. Agilidade Operacional na Transferência ou Desligamento}

Caso a colaboradora mude de função organizacional, basta revogar o papel anterior sem necessidade de auditar tabelas individuais:

\begin{lstlisting}
-- Inspecao dos privilegios operantes no papel ativo
SHOW GRANTS FOR 'sec_carla'@'localhost' USING 'r_secretaria';

-- Revogacao imediata do papel: todos os privilegios sao cancelados instantaneamente
REVOKE 'r_secretaria' FROM 'sec_carla'@'localhost';
\end{lstlisting}

\newpage

\section*{7. DESAFIO PRÁTICO: MATRIZ DE SEGURANÇA ACADÊMICA}

Para consolidar as competências de controle de acesso e segregação funcional, apresentamos o desafio que integra as tabelas \texttt{aluno}, \texttt{turma}, \texttt{matricula} e \texttt{nota}.

\subsection*{7.1. Definição da Matriz de Privilégios Corporativa}

\begin{table}[h!]
\centering
\small
\begin{tabular}{|l|l|l|}
\hline
\textbf{Papel (Role)} & \textbf{Tabela Alvo} & \textbf{Privilégios Autorizados} \\ \hline
\texttt{r\_secretaria} & \texttt{aluno, matricula} & \texttt{SELECT, INSERT, UPDATE, DELETE} (Completo) \\ \hline
\texttt{r\_secretaria} & \texttt{turma} & \texttt{SELECT, INSERT, UPDATE} (Sem exclusão física) \\ \hline
\texttt{r\_professor} & \texttt{aluno} & \texttt{SELECT} (Consulta cadastral) \\ \hline
\texttt{r\_professor} & \texttt{nota} & \texttt{SELECT, INSERT, UPDATE} (Lançamento e correção) \\ \hline
\texttt{r\_aluno} & \texttt{nota} & \texttt{SELECT} (Consulta de rendimento) \\ \hline
\end{tabular}
\caption{Matriz de controle de acesso para o subsistema acadêmico.}
\end{table}

\subsection*{7.2. Implementação dos Scripts de Concessão}

\begin{lstlisting}
CREATE TABLE nota (
    id_aluno INT,
    id_turma INT,
    valor DECIMAL(4,1) CHECK (valor BETWEEN 0 AND 10),
    PRIMARY KEY (id_aluno, id_turma)
);

-- Atribuicao conforme a matriz de seguranca
-- 1. Perfil Secretaria
GRANT SELECT, INSERT, UPDATE, DELETE ON escola.aluno TO 'r_secretaria';
GRANT SELECT, INSERT, UPDATE, DELETE ON escola.matricula TO 'r_secretaria';
GRANT SELECT, INSERT, UPDATE ON escola.turma TO 'r_secretaria';

-- 2. Perfil Docente
GRANT SELECT ON escola.aluno TO 'r_professor';
GRANT SELECT, INSERT, UPDATE ON escola.nota TO 'r_professor';

-- 3. Perfil Discente
GRANT SELECT ON escola.nota TO 'r_aluno';
\end{lstlisting}

\newpage

\section*{8. SÍNTESE DO CAPÍTULO E COMPARATIVO DE MODELOS}

\begin{table}[h!]
\centering
\small
\begin{tabular}{|l|l|p{4.5cm}|l|}
\hline
\textbf{Modelo} & \textbf{Autoridade Decisória} & \textbf{Vantagem Principal} & \textbf{Comandos Típicos} \\ \hline
\textbf{DAC} & Proprietário do Objeto & Flexibilidade de compartilhamento & \texttt{GRANT, REVOKE} \\ \hline
\textbf{MAC} & Política Central / Rótulo & Alta segurança governamental/militar & Visões, Classificação \\ \hline
\textbf{RBAC} & Função Organizacional & Escalabilidade e facilidade de gestão & \texttt{CREATE ROLE, GRANT} \\ \hline
\end{tabular}
\caption{Comparação funcional entre os modelos de controle de acesso.}
\end{table}

\section*{REFERÊNCIAS DO CAPÍTULO}

\noindent ELMASRI, R.; NAVATHE, S. B. \textbf{Sistemas de Banco de Dados}. 6. ed. São Paulo: Pearson Addison Wesley, 2011.

\vspace{0.8em}
\noindent OWASP. \textbf{OWASP Top 10: 2021 - The Ten Most Critical Web Application Security Risks}. Disponível em: \url{https://owasp.org/Top10/}. Acesso em: 13 out. 2026.

\vspace{0.8em}
\noindent SALTZER, J. H.; SCHROEDER, M. D. The protection of information in computer systems. \textbf{Proceedings of the IEEE}, v. 63, n. 9, p. 1278-1308, 1975.

\vspace{0.8em}
\noindent SANDHU, R. S. et al. Role-based access control models. \textbf{IEEE Computer}, v. 29, n. 2, p. 38-47, 1996.

\vspace{0.8em}
\noindent SILBERSCHATZ, A.; KORTH, H. F.; SUDARSHAN, S. \textbf{Sistema de Banco de Dados}. 7. ed. Rio de Janeiro: GEN LTC, 2020.

\vspace{0.8em}
\noindent STALLINGS, W.; BROWN, L. \textbf{Segurança de Computadores: Princípios e Práticas}. 4. ed. Rio de Janeiro: GEN LTC, 2018.
